Jupyter notebooks are widely used to share computational research, but a
large share of published notebooks fail to re-execute in a different
environment. The FAIR Jupyter project has studied this problem at scale and
built a pipeline that containerises each repository before execution to
remove many such failures; even so, this pipeline only detects and logs the
dependency-related failures that remain -- it does not explain or repair
them. This thesis adds a repair layer that attaches at exactly that point.

The working dataset is a residual set of 214 dependency-related failures
recorded by the Docker-based pipeline, of which 200 fall inside a pip-only
repair scope defined for this first version of the system. Given a failed
notebook, a deterministic classifier decides its subtype and repair
eligibility before any language model is consulted.

% TODO-FULL-RUN: update after the explanation component has been run over the complete dataset
\fullrunupdate{An explanation component is designed to produce a
plain-language, schema-validated description of the failure for every
classified record, independent of whether it can later be repaired; it has
not yet been run over the complete dataset.}

A repair component proposes a concrete
fix using an open,
locally-run large language model whose suggestion is grounded in live
metadata from the Python Package Index and validated, field by field,
against that evidence, so that a proposed package name or version is never
invented. A fix-application component applies a validated proposal inside
an environment rebuilt to match the one the notebook originally failed in,
re-executes the whole notebook, and classifies the result as fixed, still
failing, or not meaningfully testable.

All of these components are implemented and covered by 332 automated tests.

% TODO-FULL-RUN: update after full orchestrated dataset execution
\fullrunupdate{Each has additionally been checked in a small pilot against
real external systems: the live PyPI registry, a real local language model,
a real Docker daemon, and real GitHub repositories. This pilot shows that
each component behaves correctly against its relevant real external system,
and surfaces a recurring pattern in the dataset: a targeted repair reliably
removes the specific error it addresses, but a notebook that depends on
several packages often then reveals a second, previously hidden dependency
problem further along. Neither the pilot's scope nor this pattern has yet
been confirmed across the complete dataset.}

A result-logging component now joins the classification, explanation, and
repair records these components already produce into one persistent record
per repair attempt, stored in a SQLite benchmark table, and an RDF mapping
re-expresses that record and links it to the existing FAIR Jupyter knowledge
graph.

% TODO-FULL-RUN: update after full repair_attempts and knowledge-graph population
\fullrunupdate{Both were executed against real, pilot-scale data and
checked directly: the mapping produces syntactically valid RDF that links
correctly to existing notebook resources, but neither has yet been
populated from a full run.}

All 214 dependency-error
records were confirmed to share notebook identifiers with the inspected
knowledge-graph source data.

% TODO-FULL-RUN: rewrite once full orchestration and the full evaluation run exist
\fullrunupdate{This thesis reports work completed through
this implementation stage. Full orchestration of the pipeline over the
complete dataset, population of the benchmark table and knowledge graph
from a full run rather than a pilot sample, and a full-scale evaluation of
repair success remain future work and are stated as such, rather than
presented as already measured.}
